\section{The marked rank-three deformation ring and its equations}
\label{sec:marked-native-ring}

% BEGIN NMD1
\paragraph{NMD1. Representing the unrestricted lifts on the host quotient.}
\label{claim:NMD1}
Fix the host of CVI1--2, its coefficient DVR $O$, residue field $k$,
profinite quotient $\Gamma$, frame $F$ and a marked $g_1$. Put
$c_*=F.c$, $\eta=F.\eta$, $\Psi=F.\Psi$, and
$(w_0,w_1,w_2)=(1,0,-1)$. Write
$T(t)=\operatorname{diag}(t,1,t^{-1})$ and let $J$ reverse the three
basis vectors. The source group is
$\mathrm{GL}_3(A)\rtimes\{\pm1\}$ with the negative sign acting by
$h\mapsto J h^{-t}J$. The base representation is
$r_z(g)=(T(\Psi(g)),\eta(g))$. The target category consists of the
complete Noetherian local $O$-algebras $A$ with the given residue
identification with $k$. Its maps are local and continuous by FP1.
A lift throughout has sign $\eta$, matrices congruent to $r_z$,
and locally constant reductions at every maximal-ideal power.

Here is a finite power-series construction representing these lifts
before imposing the local and marked equations. Let $H$ be the
kernel of $\bar r_z$ including its sign. It is open normal. Write
$P$ for the intersection of kernels of all continuous maps of $H$ to
finite $\ell$-groups. This subgroup is closed and characteristic in
$H$, hence normal in $\Gamma$. Every lift kills $P$. Indeed it maps
$H$ into the congruence group, whose successive quotients are
subgroups of the additive matrix spaces over
$\mathfrak m_A^i/\mathfrak m_A^{i+1}$. These quotients are finite
$\ell$-groups, and their inverse limit is the congruence group.

The pro-$\ell$ group $H/P$ is topologically finitely generated.
CVI2 gives finite $\operatorname{Hom}_{\rm cts}(H,\mathbb F_\ell)$,
also the continuous dual of $H/P$ modulo the closed subgroup of
commutators and $\ell$th powers. This elementary abelian profinite
quotient is finite since its characters separate its points. Lift a
basis. If its lifts generated a proper closed subgroup, a finite
$\ell$-group quotient would still see it as proper; a maximal
subgroup of that finite group has index $\ell$, contradicting the
basis choice. Adjoining finitely many coset representatives gives
topological generators $g_i$ of $\Gamma/P$.

Over $B=O[[X_{iab}]]$ assign to $g_i$ the matrix
$r_z(g_i)^{\rm mat}+(X_{iab})_{ab}$ and sign $\eta(g_i)$.
The determinant is a unit by reduction. The group over $B$ is
profinite since $B$ is the inverse limit of its finite quotients.
These elements give a continuous representation of the free profinite
group on the generators. For every element of the kernel of its
surjection onto $\Gamma/P$, impose all entries of the matrix
minus the identity. Its sign already equals one. The resulting ideal
is contained in the kernel of evaluation at $X=0$ and is closed
by FP1. The tautological representation over the quotient $B_0$
therefore factors continuously through $\Gamma/P$. Its residue and
sign agree with the given ones because they agree on generators.
Every original lift factors through this group and evaluates the
variables at its generator-entry differences. It annihilates precisely
the required relators in the sense that the evaluation factors through
$B_0$. The variable values determine the algebra map uniquely.
This proves representability on the original category, including
ramified coefficient DVRs.
% END NMD1

% BEGIN NMD2
\paragraph{NMD2. Cutting out the marked normalization and all local conditions.}
\label{claim:NMD2}
For $v\in S$ use $D_v,I_v,\mathrm{loc}_v$ from CVI2 and write
$N_v=I_v\cap\ker\eta\cap\ker\bar\Psi$, with local characters
pulled back by $\mathrm{loc}_v$. At $v_\ell$ a basis congruent to
the identity must make the matrices upper triangular on $\ker\eta$,
with diagonal $\Psi(s)^{w_a}$ on $I_v\cap\ker\eta$. At
$v\in S_b$ the same basis must make them diagonal, with the same
inertia prescription. Away from $T_b=\{v_\ell\}\cup S_b$ the
representation must kill $N_v$.

We explain how the existential local bases give equations. At every
$v\in T_b$, CVI2 supplies $x\in D_v\cap\ker\eta$ with
$\Psi(x)^2-1$ a unit. The residual eigenvalues of the tautological
matrix $X=r(x)^{\rm mat}$ are distinct: $p-1,1-p^{-1}$ and
$p-p^{-1}$ are units for $p=\Psi(x)$. Hensel factorization gives
three ordered roots $\lambda_i$ of its characteristic polynomial,
with unit gaps. The idempotents
\[
 E_i=\prod_{b\ne i}(X-\lambda_b)/(\lambda_i-\lambda_b)
\]
project onto free rank-one summands reducing to $k e_i$.
Normalize $E_i e_i$ by its $i$th coordinate. Its columns form a
matrix $h\equiv 1\pmod {\mathfrak m}$ with diagonal entries one.
Uniqueness of the ordered roots and this eigenvector normalization
make $h$ commute with every target-ring map.

For any upper triangularizing basis congruent to one, its initial
coordinate flags are $X$-stable and have the prescribed residual
flags. Each such flag is the sum of its $E_i$-summands. The
summands with the required residual indices span the flag by
Nakayama; summands with other residues vanish there by the same
argument, since the flag is a direct summand. Thus these are exactly
the flags in the basis $h$. The change between the two bases is
upper triangular, preserving all diagonal entries on $\ker\eta$.
For a diagonalizing basis its individual lines agree with the
ordered eigenlines as well. Hence the original tests can all be
made in $h$, also after arbitrary coefficient base change.
Impose the lower entries on $D_v\cap\ker\eta$, also the upper
entries for $v\in S_b$, and the specified diagonal differences on
$I_v\cap\ker\eta$. On $N_v$ outside $T_b$ impose the entries
of $r(s)^{\rm mat}-1$; the multiplier there is already one.

Finally let $g_* = c_*g_1^{-1}c_*$. Impose $r(c_*)^{\rm mat}=1$
and the two entries $(0,1),(1,0)$ of
$r(g_1)^{\rm mat}-r(g_*)^{\rm mat}$ equal to zero.
These are exactly the marked normalization equations.
All the ideals just described lie in the kernel of the map to $O$
classifying $r_z$: away from $T_b$ this uses $\Psi(N_v)=1$ from
CVI2. Their sum is closed. Its quotient $R$ of $B_0$ has the same
residue field and an $O$-valued base point. For every test ring the
unique map of NMD1 factors through $R$ exactly when its representation
satisfies the marked and local conditions. This constructs the native
universal pair $(R,u)$, representing actual matrices in the given
global basis, for the same $F$ and $g_1$.
% END NMD2

% BEGIN NMD3
\paragraph{NMD3. Minimal matrix-entry coordinates and the integral origin.}
\label{claim:NMD3}
For $g\in\ker\eta$, put
$z_{gab}=u(g)^{\rm mat}_{ab}-r_z(g)^{\rm mat}_{ab}$ and give it
weight $w_a-w_b$. The $z_{gab}$ distinguish all maps out of $R$:
they determine the matrices on $\ker\eta$, and the normalized
value at $c_*$ determines them on the other coset. In particular
they span the relative cotangent space
$\mathfrak m_R/(\mathfrak m_OR+\mathfrak m_R^2)$, since a nonzero
linear functional on their quotient would give two distinct maps
to $k[\epsilon]/(\epsilon^2)$ agreeing on them (FP1--2).
Choose a basis from these entries. By complete Nakayama their
lifts give a surjection $P=O[[X_i]]\twoheadrightarrow R$ with
$X_i\mapsto z_i$. Put $J_R=\ker(P\to R)$. Each relation has linear
coefficients in $\mathfrak m_O$ by independence of the chosen
cotangent classes.

Conjugation by $T(t)$, for every unit $t$ in a test ring, multiplies
an entry of weight $j$ by $t^j$: the sign involution fixes $T(t)$,
and its diagonal form gives this formula on both cosets. This
preserves the marked equations. At an exceptional place replace a
basis $h$ by $T(t)hT(t)^{-1}$, still congruent to one, to preserve
the local equations. The outside inertia condition is preserved by
conjugation. Thus this operation acts on the represented functor,
and the chosen coordinates transform with their declared weights.

For clarity, the weight space in the source is the simultaneous test
over all ring maps whose representations differ by this conjugation.
It equals the image of the pure-weight power series in $P$.
Indeed FP3 extracts weights by the formal unit $t=1+T$.
Modulo each maximal-ideal power only finitely many monomial weights
occur. The corresponding Laurent monomials in $1+T$ are linearly
independent over the target coefficient ring: clear negative
powers and substitute $T=t-1$ in polynomials. Hence if $f\in J_R$
every weight component of $f$ vanishes in $R/\mathfrak m_R^n$
for every $n$, and therefore in $R$. If the image of $f$ has
the intrinsic weight $j$, exactly the same test shows that
its component of weight $j$ already represents it. Conversely
every pure-weight series transforms correctly, by continuity.
In particular \emph{every} matrix entry of $u$ has the prescribed
weight. Finally $r_z$ gives a map $R\to O$ killing each chosen
entry difference. Its composite with $P\to R$ is constant-term
evaluation, so every relation has zero constant term and the
augmentation sends the whole universal representation to $r_z$.
This proves the native adapted minimal presentation with origin.
% END NMD3

% BEGIN NMD4
\paragraph{NMD4. Normalized cocycles and the residual tangent inequality.}
\label{claim:NMD4}
Let $M_j$ consist of the matrices supported in weight $j$ with the
residual adjoint action. Let $H^1_L(j)$ be the native Selmer
quotient of CVI4, and $d_j=\#\{i:w(X_i)=j\}$. Then
$\dim_k H^1_L(j)\le d_j$. To prove it put $A=k[\epsilon]/(\epsilon^2)$.
A continuous cocycle $z$ gives the lift
$(1+\epsilon z(g))\bar r_z(g)$, interpreting the factor on the left
as having positive sign. The cocycle identity is exactly
multiplicativity. Replacing $z$ by a coboundary translate amounts to
conjugation by a matrix congruent to one.

Write $Z_{L,v}(j)$ for the local cocycles of CVI4. At
$v\notin T_b$ they vanish on $N_v$, preserving the outside
condition. At $v_\ell$ and $j<0$ they are coboundaries, so a
change of basis realizes the base point. At $j>0$ the correction is
strictly upper triangular, preserving the diagonal; at $j=0$ it is
diagonal and zero on inertia. For $v\in S_b$ and $j\ne0$, CVI3
gives $H^1(D_v,M_j)=0$, so the correction is a local change of
basis. For $v\in S_b$ and $j=0$ the correction is diagonal. On
inertia the action is trivial since $\eta=1$. The correction there
factors through $I_v/N_v$, of order prime to $\ell$ by CVI4, and
therefore vanishes. This verifies the exact diagonal inertia
values at these primes too. Local coboundaries preserve each test.

The following elementary normalization will also be used across
small extensions. Suppose entries are already normalized modulo a
square-zero ideal annihilated by the maximal ideal. Let $E$ be the
error in the matrix at $c_*$. Writing $Y^*=JY^tJ$, the identity
$r(c_*)^2=1$ gives $E=E^*$. Conjugation by $1+Y$ changes $E$ by
$Y+Y^*$. First take $Y=-E/2$. To preserve the resulting equality
at $c_*$ one may still use $Y=-Y^*$. On the difference of matrices
at $g_1,g_*$ this adds the commutator of $Y$ with the residual
diagonal
\[
  D=\operatorname{diag}(p-p^{-1},0,p^{-1}-p),
       \qquad p=\bar\Psi(g_1).
\]
Both elements have positive sign. The multipliers $D_{11}-D_{00}$
and $D_{00}-D_{11}$ are units in $k$. Choose $Y_{01},Y_{10}$
to cancel the two remaining errors, set their reflected entries to
their negatives, and set the other entries to zero. The two
reflected pairs are disjoint. This correction preserves weight
$j$ whenever the errors have that weight, since $(a,b)$ and
$(2-b,2-a)$ have the same weight. This procedure on residual
cocycles is $k$-linear.

Consequently the subspace of normalized Selmer cocycles surjects
onto $H^1_L(j)$. Its evaluation on the selected entries of weight
$j$ is injective. Indeed the pushed-forward $z_i$ is
$\epsilon z(g_i)_{a_i b_i}\bar\Psi(g_i)^{w_{b_i}}$, automatically
zero for all other weights. If the indicated evaluations vanish,
every coordinate has the base value, so the classifying map equals
the base map by NMD3. This forces $z=0$. Taking dimensions proves
the inequality without assuming uniqueness of a normalized
representative in a cohomology class.
% END NMD4

% BEGIN NMD5
\paragraph{NMD5. Genuine graded local lifts across a small extension.}
\label{claim:NMD5}
Here $A\to A'$ is a surjection of finite test rings, $A$ is graded
by $\mathbb Z$, and the kernel $L$ is homogeneous of a single weight
$j$, contained in $\mathfrak m_A$ and killed by $\mathfrak m_A$.
Thus $L$ is a $k$-space with $L^2=0$. Suppose a native point over
$A'$ has local matrix entries admitting lifts of their prescribed
weights. At each place there is a genuine continuous representation
over $A$ lifting its restriction, with graded entries and the
prescribed local condition.

Outside $T_b$ the given representation factors through
$D_v/N_v$. The group $I_v/N_v$ is finite of order prime to
$\ell$, and $D_v/I_v\simeq\widehat{\mathbb Z}$ on the actual host.
Hence $H^2(D_v/N_v,M_j)=0$: averaging removes positive-degree
cohomology of the finite kernel, and the procyclic quotient has
$\ell$-cohomological dimension one. For the latter fact in degree
two, a continuous extension by a finite discrete module is
profinite and a lift of the generator gives a section from the
free profinite group on one generator. Equivalently this kills
the class of every finite-valued continuous factor set.
Choose lifts of the matrices of the finite image, keeping each
entry in its indicated degree and keeping the sign. They are
invertible by reduction. The multiplication error has the form
$1+C(s,t)$ with $C$ in $M_j\otimes_k L$ (here $L$ is the scalar
multiplicity space); a nonzero kernel entry can only be in matrix
weight $j$. Associativity makes it a continuous two-cocycle for
the residual adjoint action. The vanishing just proved allows
a continuous one-cochain correction, with the same support, which
makes the lift multiplicative on $D_v/N_v$. Products, inverses
and the negative-sign involution preserve the entry degrees.
For inverses this follows by the adjugate formula, whose determinant
has degree zero and is a unit in the degree-zero subring
(the degree-zero component of its inverse is its inverse).

At $v\in T_b$ first use the element $x$ of NMD2. Lift its matrix
in the prescribed degrees. Its characteristic coefficients have
degree zero. The degree-zero subring of the finite graded local
ring is local with nilpotent maximal ideal and has residue field
$k$: an element of nonzero residue is invertible in the subring
by the preceding observation. Thus its three Hensel roots have
degree zero. The eigenvector construction in NMD2 gives a graded
matrix $h$ reducing to one modulo $\mathfrak m_A$ and lifting the
canonical eigenbasis downstairs. Indeed $(E_i)_{ai}$ has weight
$w_a-w_i$ and its diagonal normalizing factor has degree zero.

In this basis all matrices downstairs are upper triangular at
$v_\ell$, and diagonal at $S_b$, even on the negative-sign coset.
For that coset put $y=gxg^{-1}$, of positive sign. The matrix at
$g$ intertwines the diagonal $Jr(x)^{-t}J$ with the matrix at $y$.
Both have the same ordered residual eigenvalues. The matrix at $y$
is upper triangular, so its ordered eigenflag is the standard one,
as in NMD2. The intertwining therefore preserves this flag. In the
diagonal case it preserves its three ordered lines. This proves the
assertion on the extra coset.

The diagonal projection on the upper triangular semidirect group
is a homomorphism, since $h\mapsto Jh^{-t}J$ preserves the Borel
subgroup and on the diagonal reverses and inverts the entries.
Divide off the torus matrix $T(\Psi)$ from these diagonal values.
This gives a crossed homomorphism into the diagonal units with
sign action by $\eta$. It is one on $I_v$, where $\eta=1$ and the
three diagonal values are prescribed. Together with its sign it
is therefore a homomorphism on $D_v/I_v$ into diagonal units
semidirect signs. Lift its Frobenius value in the degree-zero
diagonal units, retaining the sign. This gives a continuous
homomorphism on $\widehat{\mathbb Z}$ since the target is finite.
Inflate and multiply back by $T(\Psi)$, which commutes with all
diagonal factors and is fixed by the sign involution. This lifts
the diagonal representation and all its inertia values.

For $S_b$ conjugate back by $h$ to finish. At $v_\ell$ choose
graded upper triangular set lifts with \emph{these} diagonal entries,
locally constantly by choosing once for each pair consisting of
a downstairs matrix and its lifted diagonal value. Their
multiplication error is now strictly upper triangular as well as
of weight $j$. For $j\le0$ it is zero. For $j>0$, CVI3 gives
$H^2(D_v,M_j)=0$, hence also vanishing for finitely many copies
indexed by a basis of $L$. Correct the error by a continuous
one-cochain in that strictly upper weight. This preserves the
diagonals, the inertia conditions and the lift. Conjugation back
by $h$ gives the required local representation on the actual frame.
% END NMD5

% BEGIN NMD6
\paragraph{NMD6. Detecting the equation space by the native obstruction group.}
\label{claim:NMD6}
Write $J_R=\ker(P\to R)$, $\mathfrak m=\mathfrak m_P$. By
Artin--Rees choose $n\ge1$ with $\mathfrak m^n\cap J_R\subset
\mathfrak mJ_R$. In $B=P/\mathfrak m^n$ let $J_B$ be the
image of $J_R$. The natural map
$J_R/\mathfrak mJ_R\to J_B/\mathfrak m_B J_B$ is an isomorphism.
The ring $B$ is finite, with a grading by monomial weight, as only
monomials of total degree below $n$ survive; its maximal ideal and
$J_B$ are homogeneous by NMD3. Thus the right side is a finite
graded $k$-space. For a fixed weight $j$, let $L$ be its degree-$j$
summand. In $B/\mathfrak m_BJ_B$ all the other summands together
form an ideal since multiplication acts on them through $k$.
Quotient by this ideal to obtain $A$. Put $A'=B/J_B$. The map
$A\to A'$ has kernel $L$, square zero and killed by the maximal
ideal, and the map $P\to A$ sends $J_R$ onto $L$.

Push $u$ forward along $R\to A'$ to obtain $r'$. All its entries
have lifts of the prescribed degrees by NMD3. Fix a continuous set
lift $s$ to the semidirect group over $A$ with these degrees and
the original sign, choosing once per element of the finite image.
At each place take the genuine local lift $r_v$ of NMD5. Define
$C(g,h)$ and $B_v(d)$ by the identities
\[
 s(g)s(h)s(gh)^{-1}=1+C(g,h),\qquad
 s(\mathrm{loc}_v d)r_v(d)^{-1}=1+B_v(d).
\]
The errors have sign one, lie in $M_j\otimes_k L$, and are locally
constant. Conjugation on such errors is the residual adjoint action
because $\mathfrak m_A L=0$. Associativity gives $dC=0$, and
multiplicativity of $r_v$ gives $dB_v=\mathrm{res}_v C$.
Applying a functional $\lambda:L\to k$ gives a linear map
$L^*\to F.\mathrm{Ob}(j)$ by CVI4's cochain definition.

If $0\ne\lambda$ mapped to zero, quotient $A$ by $\ker\lambda$,
which is an ideal inside the small kernel. Denote the resulting ring
by $A_\lambda$ and identify its nonzero kernel over $A'$ with $k$
via $\lambda$. The coboundary criterion of CVI4 gives a continuous
global one-cochain $h$ with $dh=\lambda C$, such that each
$\lambda B_v-\mathrm{res}_v h$ is in $Z_{L,v}(j)$. Multiplying
the pushed-forward $s(g)$ on the left by $1-h(g)$ gives a genuine
global lift. Relative to each pushed-forward $r_v$ its difference
is the indicated local cocycle in the small kernel.

These changes preserve the local conditions also over $A_\lambda$.
A small-kernel coboundary conjugates by a matrix congruent to one.
Off $T_b$ the correction vanishes on $N_v$. At $S_b$ the nonzero
weights are coboundaries by CVI3. In weight zero the correction is
diagonal and vanishes on inertia as in NMD4. At $v_\ell$ a negative
weight is a coboundary, a positive weight is strictly upper, and
weight zero is diagonal and zero on inertia. In the local triangular
or diagonal basis the error matrix is unchanged: that basis reduces
to the identity, and its maximal-ideal entries annihilate the
kernel. Thus the fixed inertial diagonals and the required local
bases survive. Apply the normalization in NMD4 at $c_*,g_1$;
it preserves the local conditions and reduces to the identity over
$A'$. We now have a native point over $A_\lambda$ lifting $r'$.

By NMD2 this point gives an algebra map $R\to A_\lambda$. Its
composite with $P\to R$ and the original quotient map
$P\to A_\lambda$ agree after projection to $A'$ by uniqueness
for $r'$. The difference of their values on any series $f$ is
the sum of its linear coefficients times the corresponding
coordinate differences: constants agree, and all terms of degree
at least two have zero difference because each coordinate is in
the maximal ideal and each difference lies in the small kernel.
Equivalently both maps agree on $\mathfrak m^2$; they are
continuous. By minimality in NMD3 they therefore agree on $J_R$.
One map kills $J_R$, while the original quotient sends it onto
$L/\ker\lambda\ne0$. This contradiction proves
\[
 \dim_k (J_R/\mathfrak mJ_R)_j\le\dim_k F.\mathrm{Ob}(j).
\]
% END NMD6

% BEGIN NMD7
\paragraph{NMD7. The five multiplicities for the same presentation.}
\label{claim:NMD7}
Choose a homogeneous basis of each summand of $J_R/\mathfrak mJ_R$
constructed in NMD6. Each basis element lifts to a series in $J_R$
of its precise weight: start with any lift and take its
weight component as in NMD3. The finite graded truncation used in
NMD6 commutes with that projection. By Nakayama these finitely many
lifts generate $J_R$ as an ideal, since $P$ is Noetherian.
Let $b_j=\dim_k F.\mathrm{Ob}(j)$. Pad these lists with zero
series to obtain exactly $b_j$ homogeneous relations at each weight
$j$. Their constant terms are zero by NMD3. Their linear
coefficients are in $\mathfrak m_O$ by minimality.

Every variable is an entry difference with weight in $[-2,2]$.
The obstruction space outside this interval is zero since all
cochains there take values in the zero matrix weight space.
By CVI4 the dimension of $H^1_L(j)$ equals $b_j-1$ for $j=-1$,
$b_j+1$ for $j=1,2$, and $b_j$ otherwise. Combining with NMD4
gives the five bounds on the numbers of generating equations:
\[
 b_{-2}\le d_{-2},\quad b_{-1}\le d_{-1}+1,\quad
 b_0\le d_0,\quad b_1\le d_1-1,\quad b_2\le d_2-1.
\]
In particular $d_1,d_2\ge1$. The marked pair, origin, homogeneous
equations and residual tangent bounds needed in CVI5 now all belong
to one algebra on the original arithmetic frame. The nonpositive
\emph{integral} tangent test uses in addition the displayed
$F.\mathrm{VanE}$ input; it is proved in VSCM4.
% END NMD7
